% !TEX root = ../bb_AiRa_Kappaleet.tex

%% HARRIS: Chapter 10

% ö = \%c3\%b6
% ä = \%C3\%A4

\xdef\sectiontitle{\IIsectiontitle} %aiheuttama
\TOCrefs

%%%%%%%%%%%%%%%
\begin{frame}[label=2_9_a]%[label=2_8_TOC1]
\frametitle{\texorpdfstring{\culine[\sectiontitleulinecolor]{\sectiontitle}}{\sectiontitle} }
\label{\TOCname}

\vspace{\chapterTOCvksip}
\begin{itemize}[leftmargin=0.5cm,itemsep=3mm]
    \item[2.1]<1-> \hyperlink{2_1_a}{\IIaaSubsectiontitle}
    \item[2.2]<1-> \hyperlink{2_2_a}{\IIaSubsectiontitle}
    \item[2.3]<1-> \hyperlink{2_3_a}{\IIbSubsectiontitle}	
    \item[2.4]<1-> \hyperlink{2_4_a}{\IIcSubsectiontitle}	
    \item[2.5]<1-> \hyperlink{2_5_a}{\IIdSubsectiontitle}
    \item[2.6]<1-> \hyperlink{2_6_a}{\IIeSubsectiontitle}
    \item[2.7]<1-> \hyperlink{2_7_a}{\IIfSubsectiontitle}
    \item[2.8]<1-> \hyperlink{2_8_a}{\IIgSubsectiontitle}
    \item[2.9]<1-> \hyperlink{2_9_a}{\CDAlert[red]{\IIhSubsectiontitle}}
    \item[2.10]<1-> \hyperlink{2_10_a}{\IIiSubsectiontitle}
\end{itemize}

%\endofslide 
\end{frame}
%%%%%%%%%%%%%%%

\xdef\sectiontitle{\IIhSubsectiontitle} 

%%%%%%%%%%%%%%%
\begin{frame}
\frametitle{\texorpdfstring{\culine[\sectiontitleulinecolor]{\sectiontitle}}{\sectiontitle} }
\framesubtitle{Suprajohtavat alkuaineet}%Superconducting elements
%Superconducting elements
\appear{}
\begin{itemize}
	\item Joidenkin metallisten alkuaineiden \appear{resistiivisyys voi yllättäen laskea nollaan kun lämpötila lähestyy absoluuttista nollapistettä $T=0$\;K}% 
	\appear{\xdef\loc{south east}
\begin{tikzpicture}[remember picture,overlay] % anchor=south
    \node (A) [yshift=-0.25*\figYshift,xshift=\figXshift, anchor=\loc] at (current page.\loc) {\includegraphics[page=1,width=6.7cm]{Figures_booklet/SOM_10_Bonding_figures/SOM_Bonding_Figure_48.png}}; %scale=\figscale. % 8cm max height
\end{tikzpicture}%
\footnote{\HarrisCR[1-]}}%
\appear{(vertaa tinaa ja kuparia) }
%	The critical temperature for this to happen, $T_{c}$, is a characteristic of the material. For those elements which behave like this, the critical temperature $T_{c}$ is usually in the order of few Kelvins. $0.15$
\end{itemize}

\vspace{2.5mm}
\begin{minipage}{7cm}
\begin{itemize}%[itemsep=1.7mm]
\item  \CDAlert[red]{Kriittinen lämpötila} \appear{(\CDAlert[red]{$T_{c}$})} \appear{ilmiön ta\-pah\-tu\-miselle} \appear{ riippuu materiaalin omi\-nai\-suuk\-sista}

\item Näin käyttäytyvillä alkuaineilla\appear{ kriit\-tinen lämpötila $T_{c}$ on yleensä \\ \CDAlert[red]{muutaman kelvinin luokkaa}}
	\end{itemize}
\end{minipage}

\endofslide 
%\endofsection
\end{frame}
%%%%%%%%%%%%%%

%%%%%%%%%%%%%%%%
\begin{frame}
\frametitle{\texorpdfstring{\culine[\sectiontitleulinecolor]{\sectiontitle}}{\sectiontitle} }
\framesubtitle{Suprajohtavat alkuaineet}
\appear{}
\begin{itemize}[itemsep=1.9mm]
	\item Yli 20 \CDAlert[red]{metallista alkuainetta} voi muuttua suprajohteeksi (superconductor, SC) \appear{(taulukko)%
	\xdef\loc{south east}
\begin{tikzpicture}[remember picture,overlay] % anchor=south
    \node (A) [yshift=-0.25*\figYshift,xshift=\figXshift, anchor=\loc] at (current page.\loc) {%
    {\scriptsize
\begin{tabular}{llr} 
Tyyppi & Materiaali & $T_{c}(\mathrm{~K})$ \\
\hline Yksinkertainen & $\mathrm{Al}$ & $1,17$ \\
metallinen & $\mathrm{ln}$ & $3,41$ \\
suprajohde &  $\mathrm{Hg}$ & $4,20$ \\
& $\mathrm{Sn}$ & $3,72$ \\
& $\mathrm{Pb}$ & $7,20$ \\
& $\mathrm{Nb}$ & $9,25$ \\
& $\mathrm{Nb}_{3} \mathrm{Sn}$ & $18,0$ \\
& $\mathrm{Nb}_{3} \mathrm{Ge}$ & $23,2$ \\
& $\mathrm{C}_{60} \mathrm{Rb}_{3}$ & 31 \\
& $\mathrm{La}_{2-\mathrm{S}} \mathrm{Sr}_{\mathrm{x}} \mathrm{CuO}_{4}$ & 38 \\
& $\mathrm{YBa}_{2} \mathrm{Cu}_{3} \mathrm{O}_{7}$ & 93 \\
& $\mathrm{Bi}_{2} \mathrm{Sr}_{2} \mathrm{Ca}_{2} \mathrm{Cu}_{3} \mathrm{O}_{10}$ & 107 \\
& $\mathrm{Tl}_{2} \mathrm{Ba}_{2} \mathrm{Ca}_{2} \mathrm{Cu}_{3} \mathrm{O}_{10}$ & 125 \\
& $\mathrm{HgBa}_{2} \mathrm{Ca}_{2} \mathrm{Cu}_{3} \mathrm{O}_{8}$ & 135 \\
Magnesiumdiboridi & $\mathrm{MgB}_{2}$ & 39 \\
Rautapohjaiset & $\mathrm{FeSe}$ & 8 \\
suprajohteet & $\mathrm{LaO}_{0,89} \mathrm{~F}_{0,11} \mathrm{FeAs}$ & 26 \\
& $\mathrm{Sr}_{0,5} \mathrm{Sm}_{0,5} \mathrm{FeAsF}$ & 56
\end{tabular}
}%
}; %scale=\figscale. % 8cm max height
\end{tikzpicture}}%
	\item Myös jotkin \CDAlert[blue]{puolijohteet} voivat olla suprajohtavia \appear{oikeissa olosuhteissa (esim. korkeissa pai-\\}
	\appear{neissa, tai ohutkalvorakenteina)}
	
\end{itemize}

\vspace{1.9mm}
\begin{minipage}{6.4cm}
\begin{itemize}[itemsep=1.9mm]
\item Myös \CDAlert[fuchsia]{2D materiaalit} \appear{(kaksi/kolme tietyssä kulmassa olevaa gra\-fee\-ni\-kalvoa)} \appear{voivat muodostaa suprajohteita (vuodelta 2018)}
	
	\item Niin sanottu \CDAlert[fuchsia]{korkean lämpötilan} \Alert[fuchsia]{suprajohtavuus} ei ole vielä hyvin ymmärretty ilmiö
	
	\end{itemize}
\end{minipage}

\endofslide 
%\endofsection
\end{frame}
%%%%%%%%%%%%%%%

%%%%%%%%%%%%%%%
\begin{frame}
\frametitle{\texorpdfstring{\culine[\sectiontitleulinecolor]{\sectiontitle}}{\sectiontitle} }
\framesubtitle{Suprajohteiden ominaisuuksia}%Properties of superconductors
\appear{}
\begin{itemize}
	\item \CDAlert[tunired]{Ominaisuus 1:} \appear{\CDAlert[red]{Ei sähköistä resistiivisyyttä}}
	\appear{\xdef\loc{south east}
\begin{tikzpicture}[remember picture,overlay] % anchor=south
    \node (A) [yshift=-0.25*\figYshift,xshift=\figXshift, anchor=\loc] at (current page.\loc) {%
    {\scriptsize
\begin{tabular}{lllr} 
Materiaali & Metalli/eriste/SC & $\rho(\mathrm{ohm} \mathrm{m})$ & $T_{\mathrm{c}}(\mathrm{K})$ \\
\hline hopea & $\mathrm{M}$ & $1,6 \times 10^{-8}$ & $-$ \\
Kupari & $\mathrm{M}$ & $1,7 \times 10^{-8}$ & $-$ \\
Kulta & $\mathrm{M}$ & $2,4 \times 10^{-8}$ & $-$ \\
Alumiini & M/SC & $2,8 \times 10^{-8}$ & $1,17$ \\
Lyijy & M/SC & $2,2 \times 10^{-7}$ & $7,20$ \\
Elophopea & M/SC & $9,8 \times 10^{-7}$ & $4,20$ \\
Posliini & Eriste & $5 \times 10^{12}$ & $-$ \\
Kumi & Eriste & $6 \times 10^{14}$ & $-$ \\
Kvartsi (sulatettu) & Eriste & $7,5 \times 10^{17}$ & $-$ \\
\hline
\end{tabular}
}%
}; %scale=\figscale. % 8cm max height
\end{tikzpicture}}

\item Puolijohde käyttäytyy kuin sillä  \Alert[red]{ei olisi mitattavaa (tasavirta)resistiivisyyttä} 

\item Suprajohteisiin on muodostettu virtoja, joiden ei ole kokeellisesti nähty juurikaan heikkenevän useankaan vuoden kuluessa

\item Tämä ominaisuus on \CDAlert[fuchsia]{teknologisesti relevantti}\appear{, koska resistiivisyys aiheuttaa (ohmisia) häviöitä ja laitteiden lämpenemistä}
\implies{ Sähkömagneetteihin perustuvat \\
  	laitteet voivat hyödyntää \\ suprajohteita } 

\end{itemize}



\endofslide 
%\endofsection
\end{frame}
%%%%%%%%%%%%%%%

%%%%%%%%%%%%%%%
\begin{frame}
\frametitle{\texorpdfstring{\culine[\sectiontitleulinecolor]{\sectiontitle}}{\sectiontitle} }
\framesubtitle{Suprajohteiden ominaisuuksia}
\appear{}
\semismall
\begin{itemize}[itemsep=1.6mm]
	\item \CDAlert[tunired]{Ominaisuus 2:} \appear{\Alert[red]{Täydellinen diamagnetismi} }
	\appear{\xdef\loc{south east}
\begin{tikzpicture}[remember picture,overlay] % anchor=south
    \node (A) [yshift=-0.35*\figYshift,xshift=\figXshift, anchor=\loc] at (current page.\loc) {\includegraphics[page=1,width=7.3cm]{Figures_booklet/SOM_10_Bonding_figures/SOM_Bonding_Figure_49.png}}; %scale=\figscale. % 8cm max height
\end{tikzpicture}
\footnote{\HarrisCR[1-]}}
	\item Kun termisessä tasapainossa olevaan näytteeseen kohdistetaan ulkoinen magneettikenttä, sen pinnalle muodostuu sähkövirtoja (silmukoita)
	
	\item Nämä virrat indusoivat magneettikentän \appear{joka täydellisesti sammuttaa ulkoisen magneettikentän suprajohteen sisältä }

\end{itemize}

\vspace{1.6mm}
\begin{minipage}{6.4cm}
\begin{itemize}[itemsep=1.6mm]
	\item Kun $T<T_{c}$ ulkoinen magneettikenttä \Alert[red]{kaareutuu kappaleen ympärille}

\item Jos $R=0$\appear{, niin Faradayn lain nojalla} \appear{äärettömän voimakkaita pintavirtoja voi muodostua}\appear{, mikä mahdollistaa sisäkenttien täydellisen sammumisen}

\item Suprajohteissa olevat virtasilmukat voivat vastustaa ulkoista kenttää \appear{koska $R \Approx 0$.}\appear{ Tätä suprajohteen täydellistä magneettikentän hylkimistä kutsutaan nimellä \CDAlert[red]{Meissnerin efektiksi}}
	\end{itemize}
\end{minipage}


\endofslide 
%\endofsection
\end{frame}
%%%%%%%%%%%%%%%

%%%%%%%%%%%%%%%
\begin{frame}
\frametitle{\texorpdfstring{\culine[\sectiontitleulinecolor]{\sectiontitle}}{\sectiontitle} }
\framesubtitle{Suprajohteiden ominaisuuksia}
\appear{}
\begin{itemize}
	\item Meisnerin efektin takia \appear{ suprajohde \CDAlert[red]{hylkii aina magneettia}}

\item Tällöin kestomagneetti jää \Alert[red]{levitoimaan} suprajohteen päälle
\end{itemize}

\vspace{2.5mm}
\begin{minipage}{6.4cm}
\begin{itemize}%[itemsep=1.6mm]
	\item Kun magneetti on suprajohteen päällä\appear{, se indusoi suprajohteeseen virtasilmukoita }
	
	\item Nämä \Alert[fuchsia]{virrat synnyttävät Meissnerin efektin} \appear{mikä synnyttää magneettiin suprajohdetta hylkivän voiman }
	
	\item Tietyllä etäisyydellä hylkivä voima on yhtäsuuri gravitaatiovoiman kanssa\appear{, ja magneetti alkaa le\-vi\-toi\-maan}
	\end{itemize}
\end{minipage}

\xdef\loc{south east}
\begin{tikzpicture}[remember picture,overlay] % anchor=south
    \node (A) [yshift=-0.35*\figYshift,xshift=\figXshift, anchor=\loc] at (current page.\loc) {\includegraphics[page=1,width=7cm]{Figures_booklet/SOM_10_Bonding_figures/SOM_Bonding_Figure_50.png}}; %scale=\figscale. % 8cm max height
\end{tikzpicture}
\footnote{\HarrisCR[1-]}

\endofslide 
%\endofsection
\end{frame}
%%%%%%%%%%%%%%%

%%%%%%%%%%%%%%%
\begin{frame}
\frametitle{\texorpdfstring{\culine[\sectiontitleulinecolor]{\sectiontitle}}{\sectiontitle} }
\framesubtitle{Ominaisuus 3} %Property number 3
%Property No3
\appear{}
\begin{itemize}
	\item Suprajohde käyttäytyy usein niin kuin sillä olisi \Alert[fuchsia]{$2 \Delta$-levyinen energiarako keskellä sen fermienergiaa} sen mahdollisissa \CDAlert[fuchsia]{yhden elektronin tiloissa} 
	
	\item Tällöin suprajohteen elektronin energiaa $\mathcal{E}$ voidaan kasvattaa  \appear{(tai vähentää)} \appear{vain jos $\mathcal{E}-E_{F}$} \appear{(tai $E_{F}-\mathcal{E}$) ylittää arvon $\Delta$ }
	\item Energiarako $\Delta$ kasvaa lämpötilan laskiessa\appear{, tasoittuen maksimiarvoonsa $\Delta(0)$ todella matalissa lämpötiloissa}
\end{itemize}

\xdef\loc{south}
\begin{tikzpicture}[remember picture,overlay] % anchor=south
    \node (A) [yshift=-0.15*\figYshift,xshift=0*\figXshift, anchor=\loc] at (current page.\loc) {\includegraphics[page=1,width=10cm]{Figures_booklet/SOM_BCS_band_gap_X_Zhang_phd_thesis.png}};%SOM_supeconductor_deltaE.png
\end{tikzpicture}
\footnote{\tiny Figure from the \href{https://scholar.colorado.edu/concern/graduate\_thesis\_or\_dissertations/2f75r8102}{PhD Thesis of X. Zhang}.}%(\EMPH{link}{https://scholar.colorado.edu/concern/graduate\_thesis\_or\_dissertations/2f75r8102})

\endofslide 
%\endofsection
\end{frame}
%%%%%%%%%%%%%%%

%%%%%%%%%%%%%%%
\begin{frame}
\frametitle{\texorpdfstring{\culine[\sectiontitleulinecolor]{\sectiontitle}}{\sectiontitle} }
%\framesubtitle{Property number 3}
\appear{}
\begin{itemize}
	\item Siirtymän kriittinen lämpötila madaltuu ulkoisen magneettikentän kasvaessa\appear{, ja lopulta kentän ylittäessä \CDAlert[red]{kriittisen arvon} $B_{c}$}\appear{ vaihe\-muu\-tosta ei enää tapahdu.} \appear{Voi ajatella niin että ulkoinen $B$-kenttä on niin vahva että se hävittää suprajohtavuuden}
\end{itemize}

\vspace{2.5mm}
\begin{minipage}{7.1cm}
\begin{itemize}%[itemsep=1.6mm]
	\item Kriittisen kenttäarvon $B_{c}$ lämpötilariippuvuus on likimain muotoa
\[
B_{c}(T)  \equal \appear{B_{c}(0)} \LP \appear{1-} \frac{T^{2}}{T_{c}^{2}} \;\RP
\]

\item Huomaa että koska materiaalit kestävät vain tietyn määrän magneettikenttää\appear{, vastaavasti on olemassa myös \CDAlert[red]{maksimivirta} joka materiaalissa voi kulkea/kiertää}

	\end{itemize}
\end{minipage}

\xdef\loc{south east}
\begin{tikzpicture}[remember picture,overlay] % anchor=south
    \node (A) [yshift=-0.35*\figYshift,xshift=\figXshift, anchor=\loc] at (current page.\loc) {\includegraphics[page=1,width=6.5cm]{Figures_booklet/SOM_Tipler_Solid_state_physics_figures/SOM_Tipler_SSP_figure_61.png}}; %scale=\figscale. % 8cm max height
\end{tikzpicture}
\footnote{\TiplerCR[1-]}

\endofslide 
%\endofsection
\end{frame}
%%%%%%%%%%%%%%%

%%%%%%%%%%%%%%%
\begin{frame}
\frametitle{\texorpdfstring{\culine[\sectiontitleulinecolor]{\sectiontitle}}{\sectiontitle} }
\framesubtitle{Tyypin-I suprajohteet}%I:n lajin suprajohteet,  Type-I superconductors
\appear{}
%Type-| Superconductors
\begin{itemize}[itemsep=1.8mm]
	\item \CDAlert[red]{Tyypin-I suprajohteilla tapahtuu täydellinen Meissnerin efekti} 
	
	\item Kun ulkoinen magneettikenttä on pienempi kuin $B_{c}$ \appear{($B<B_c$)}\appear{, suprajohteen sisälle indusoitunut kenttä on yhtäsuuri ja vastakkaismerkkinen ulkoiseen kenttään nähden.} \appear{Tyypin-I suprajohteille on tyypillistä että ulkoisen kentän kasvaessa vaihesiirtymä tapahtuu nopeasti}
\end{itemize}

\vspace{1.8mm}
\begin{minipage}{4.6cm}
\begin{itemize}[itemsep=1.8mm]
	\item Suprajohtavat alkuaineet ovat usein tyypin-I suprajohteita
	
	\item Niiden kriittiset kentät ja lämpötilat ovat suhteellisen matalia\appear{, tyypillisesti 0,01–0,1\;T ja 1–9\;K}\appear{, joten niiden \Alert[red]{käyttöarvo on rajallinen}}
\end{itemize}
\end{minipage}

\xdef\loc{south east}
\begin{tikzpicture}[remember picture,overlay] % anchor=south
    \node (A) [yshift=-0.1*\figYshift,xshift=\figXshift, anchor=\loc] at (current page.\loc) {\includegraphics[page=1,width=9cm]{Figures_booklet/SOM_10_Bonding_figures/SOM_Bonding_Figure_51.png}}; %scale=\figscale. % 8cm max height
\end{tikzpicture}
\footnote{\HarrisCR[1-]}

\endofslide 
%\endofsection
\end{frame}
%%%%%%%%%%%%%%%

%%%%%%%%%%%%%%%
\begin{frame}
\frametitle{\texorpdfstring{\culine[\sectiontitleulinecolor]{\sectiontitle}}{\sectiontitle} }
\framesubtitle{Tyypin-II suprajohteet}
%Type-II Superconductors
\appear{}
\seminormal
\begin{itemize}[itemsep=1.4mm]
	\item \Alert[red]{Kun $B<B_{c1}$ (ja $T<T_{c}$), tyypin-II aineissa voi tapahtua täydellinen Meissnerin efekti}

\item Tällöin materiaalin sisäinen kenttä on nolla ($B=0$)

\item Kun $B_{c 1}<B<B_{c 2}$ \appear{kenttäviivat alkavat lävistää materiaalin} \appear{mutta \Alert[blue]{viivat ovat rajoittuneet tavallisen aineen tavoin käyttäytyvien vuoputkien sisälle}} \appear{jotka muodostavat niin sanotun \CDAlert[blue]{pyörrehilan}} \appear{(engl. vortex lattice)}
\end{itemize}

\vspace{1.4mm}
\begin{minipage}{4.6cm}
\begin{itemize}[itemsep=1.4mm]
	\item Lämpötilassa $T<T_{c}$\appear{, kun $B$-kenttä lähestyy arvoa $B_{c 2}$} \appear{suprajohtavan aineen määrä pienenee}\appear{, ja vastaavasti vuoputkien sisältämä tilavuus kasvaa}
	 \item Kun $B>B_{c 2}$ \appear{koko materiaali käyttäytyy tavallisen aineen tavoin}
\end{itemize}
\end{minipage}

\xdef\loc{south east}
\begin{tikzpicture}[remember picture,overlay] % anchor=south
    \node (A) [yshift=-0.1*\figYshift,xshift=\figXshift, anchor=\loc] at (current page.\loc) {\includegraphics[page=1,width=8.8cm]{Figures_booklet/SOM_10_Bonding_figures/SOM_Bonding_Figure_54.png}}; %scale=\figscale. % 8cm max height
\end{tikzpicture}
\footnote{\HarrisCR[1-]}

\endofslide 
%\endofsection
\end{frame}
%%%%%%%%%%%%%%%

%%%%%%%%%%%%%%%
\begin{frame}
\frametitle{\texorpdfstring{\culine[\sectiontitleulinecolor]{\sectiontitle}}{\sectiontitle} }
\framesubtitle{Tyypin-II suprajohteet}
\appear{}
\seminormal

\vspace{1.5mm}
\begin{minipage}{9.1cm}
\begin{itemize}[itemsep=2mm]
	\item Kun magneettikenttä on välillä $B_{c 1}<B<B_{c 2}$, suprajohteen ei-suprajohtaviin alueisiin\appear{ muodostuu niin sanottuja \CDAlert[fuchsia]{normaaleja ytimiä (normal cores)}}\appear{, joiden läpi kenttäviivat kulkevat} 
	
	\item Normaalien ytimien ympärille muodostuu \CDAlert[fuchsia]{pyörretiloja (vorteksitiloja)}\appear{, joissa esiintyy ulkoista kenttää heikentäviä virtasilmukoita} 
\end{itemize}
\end{minipage}

\vspace{3mm}
\begin{minipage}{6.5cm}
\begin{itemize}[itemsep=2mm]
		\item Kriittisen arvon $B_{c 1}$ yläpuolella\appear{ kenttäviivat alkavat läpäistä materiaalin normaalien ytimien muodostuessa}\appear{, kun taas niiden \Alert[red]{ulkopuoliset alueet säilyttävät suprajohtavuutensa}}
	
	\item  Tätäkin suuremmilla B-kentän arvoilla \appear{ pyörteet tiivistyvät (kuva)}\appear{ ja suprajohtavan aineen \Alert[red]{alue pienenee}}
\end{itemize}
\end{minipage}

\xdef\loc{north east}
\begin{tikzpicture}[remember picture,overlay] % anchor=south
    \node (A) [yshift=0.7*\figYshift,xshift=1*\figXshift, anchor=\loc] at (current page.\loc) {\includegraphics[page=1,height=4.25cm]{Figures_booklet/SOM_10_Bonding_figures/SOM_Bonding_Figure_52.png}}; %scale=\figscale. % 8cm max height
\end{tikzpicture}

\xdef\loc{south east}
\begin{tikzpicture}[remember picture,overlay] % anchor=south
    \node (A) [yshift=-0.27*\figYshift,xshift=\figXshift, anchor=\loc] at (current page.\loc) {\includegraphics[page=1,height=4.4cm]{Figures_booklet/SOM_10_Bonding_figures/SOM_Bonding_Figure_53.png}}; %scale=\figscale. % 8cm max height
\end{tikzpicture}
\footnote{\HarrisCRshort[1-]}

\endofslide 
%\endofsection
\end{frame}
%%%%%%%%%%%%%%%

%%%%%%%%%%%%%%%
\begin{frame}
\frametitle{\texorpdfstring{\culine[\sectiontitleulinecolor]{\sectiontitle}}{\sectiontitle} }
\framesubtitle{BCS-teoria}
\appear{}
%BSC Theory
\begin{itemize}[itemsep=1.5mm]
	\item \Alert[fuchsia]{Ensimmäinen suprajohtavuuteen liittyvä Nobel-palkinto myönnettiin jo 1913} (Kammerlingh Onnesille) \appear{kun hän löysi joidenkin alkuaineiden suprajohtavuuden matalissa lämpötiloissa}
	
	\item Myöhemmin selvisi että matalan lämpötilan suprajohtavuuden selittää aineessa ilmenevät elektroniparit \appear{eli \CDAlert[red]{Cooperin parit}}\appear{, ja elektronien väliset korrelaatiot/vuorovaikutukset} \appear{(\Alert[fuchsia]{toinen Nobel-palkinto suprajohtavuudelle vuonna 1972})}

\item Elektronit ovat fermioneita\appear{ mutta \CDAlert[blue]{elektroniparit} \Alert[blue]{muodostavat bosoneita}!}

 \item Parin elektronikorrelaatiot voivat säilyä pitkiä matkoja 
 
 \item Ilmiö on seurausta \Alert[red]{elektronien ja fononien välisistä vuorovaikutuksista} \appear{(fononi on atomihilan ionien hilavärähtelykvantti).}\appear{ Kun elektronit vuorovaikuttavat keskenään}\appear{, niiden välille muodostuu heikko puoleensavetävä voima}
\end{itemize}

\endofslide 
%\endofsection
\end{frame}
%%%%%%%%%%%%%%%

%%%%%%%%%%%%%%%
\begin{frame}
\frametitle{\texorpdfstring{\culine[\sectiontitleulinecolor]{\sectiontitle}}{\sectiontitle} }
\framesubtitle{BCS-teoria}
\appear{}
\begin{itemize}
	\item Kaksi vihjettä opasti tutkijat ajatukseen fononien \appear{(atomihilan värähtelyiden)} \appear{välittämistä elektronien välisistä vuorovaikutuksista} 
	
	\item \CDAlert[red]{Vihje 1}: \appear{Materiaalin kriittinen lämpötila $T_c$ riippuu atomihilan ionien keskimääräisestä massasta:} $\appear{T_{c}} \propto \appear{M^{-1 / 2}}$\appear{, mikä havaittiin \CDAlert[red]{isotooppikokeilla}}

\item Tarkalleen ottaen riippuvuus on muotoa $ M^{\alpha} T_{c}=$ vakio\appear{,
missä eksponentti $\alpha$ riippuu materiaalista}
\implies{ Tämä kokeellinen havainto atomimassojen vaikutuksesta $T_{C}$:an implikoi että \\
	atomihilalla on ilmiön muodostumisessa rooli}
	
\item \CDAlert[blue]{Vihje 2}: \appear{Kokeiden perusteella matalan $T$:n suprajohtavat alkuaineet ovat \Alert[blue]{huonoja johteita huoneenlämmössä}}\appear{, ja päinvastoin} 
\implies{ Vahva vuorovaikutus elektronien ja ionien välillä huoneenlämmössä on, enemmän\\ tai vähemmän tarpeellista matalan T:n suprajohtavuudessa}
	
	\end{itemize}

\endofslide 
%\endofsection
\end{frame}
%%%%%%%%%%%%%%%

%%%%%%%%%%%%%%%
\begin{frame}
\frametitle{\texorpdfstring{\culine[\sectiontitleulinecolor]{\sectiontitle}}{\sectiontitle} }
\framesubtitle{BCS-teoria}
\appear{}
\seminormal
\begin{itemize}[itemsep=1.5mm]
	\item Matalissa $T$:ssa kaksi elektronia voi muodostaa Cooperin parin%, mediated by a phonon 
	\item Elektroni liikkuu aineen läpi \CDAlert[red]{muokaten paikallisesti atomihilaa}\appear{, mikä vaikuttaa muiden elektronien etenemiseen}
\end{itemize}

\vspace{1.9mm}
\begin{minipage}{6.7cm}
\begin{itemize}[itemsep=2.2mm]
\item Elektronit emittoivat fononeita\appear{, minkä toinen elektroni absorboi} 

\item Kokonaisuutena prosessi näyttäytyy elektroneita \Alert[blue]{puoleensavetävänä voimana} \appear{(vaikuttaa analogiselta voimien kanssa joita virtuaalisten hiukkasten vaihdot synnyttävät kvanttikenttäteorioissa)}

\item Cooperin parin \CDAlert[fuchsia]{sitoutumisenergia} on pieni \appear{(${\sim}1$\;meV)}\appear{, ja elektronien välinen etäisyys aika suuri}\appear{ (${\sim}\,1\;\mu \mathrm{m}$).} \appear{Monet Cooperin parit menevät siis limittäin.} %(see Slide 12)
\end{itemize}
\end{minipage}

\xdef\loc{south east}
\begin{tikzpicture}[remember picture,overlay] % anchor=south
    \node (A) [yshift=-0.4*\figYshift,xshift=\figXshift, anchor=\loc] at (current page.\loc) {\includegraphics[page=1,width=6.8cm]{Figures_booklet/SOM_10_Bonding_figures/SOM_Bonding_Figure_55.png}}; %scale=\figscale. % 8cm max height
\end{tikzpicture}
\footnote{\HarrisCR[1-]}

\endofslide 
%\endofsection
\end{frame}
%%%%%%%%%%%%%%%

%%%%%%%%%%%%%%%
\begin{frame}
\frametitle{\texorpdfstring{\culine[\sectiontitleulinecolor]{\sectiontitle}}{\sectiontitle} }
\framesubtitle{BCS-teoria}
\appear{}
\begin{itemize}
	\item Vuoden 1972 fysiikan Nobel-palkinto myönnettiin \Alert[red]{John Bardeenille, Leon Cooperille ja John Robert Schriefferille heidän \CDAlert[red]{BCS-teoriastaan}}\appear{, joka selitti kahden tyyppisten suprajohteiden olemassaolon}\appear{, Meissnerin efektin}\appear{, kriittisen lämpötilan olemassaolon, jne.} \appear{Alla joitakin BCS-teorian pääkohdista:}

\item[] \begin{itemize} \seminormal
	\item Systeemin globaalisti \CDAlert[red]{matalimman energian} tila tapahtuu kun \appear{energeettisimmät johtavuuselektronit (lähellä $E_{F}$:a) muodostavat \Alert[red]{vastakkaisen spinin Cooperin pareja}}
	
\item Vain \Alert[blue]{fermienergian $E_{F}$ lähellä olevat elektronit osallistuvat} suprajohtavaan tilaan. \appear{Matalampienergisten elektronien nousu lähelle fermienergiaa}\appear{ kasvattaisi energiaa enemmän kuin minkä heikko elektronien vuorovaikutus kestäisi}

\item Cooperin parin \CDAlert[fuchsia]{sitoutumisenergia} on pieni \appear{(${\sim}1$\;meV)}\appear{, ja elektronien välinen etäisyys aika suuri} \appear{(${\sim}\,1\;\mu \mathrm{m}$).} \appear{Monet Cooperin parit menevät siis limittäin.} %(see Slide 12)

\end{itemize}
\end{itemize}

\endofslide 
%\endofsection
\end{frame}
%%%%%%%%%%%%%%%

%%%%%%%%%%%%%%%
\begin{frame}
\frametitle{\texorpdfstring{\culine[\sectiontitleulinecolor]{\sectiontitle}}{\sectiontitle} }
\framesubtitle{BCS-teoria}
\appear{}
\begin{itemize}
	\item[] \begin{itemize}  \seminormal
	\item \Alert[blue]{Tyypin-I suprajohteissa olevien elektronien keskimääräiset vapaat matkan ovat suhteellisen pitkiä}\appear{, tyypin-II suprajohteissa matkat ovat lyhyempiä}\appear{, mikä mahdollistaa mikroskooppisen normaalien ytimien (B-kenttä läpäisi suprajohteen näiden kautta) verkoston muodostumisen suprajohtavan aineen sisälle}
	
	\item Cooperin parien sitoutumisenergia \appear{(eli suprajohtavuuden energiarako)} \appear{on $\Approx 3.5 k_{B} T_{c}$}\appear{, johtaen $\Approx 10^{-3} \;\mathrm{eV}$.} \appear{\Alert[fuchsia]{Sitoutumisenergian olemassaolo näkyy vasta matalissa lämpötiloissa}}
	
\item BCS-teorian mukaan \appear{\Alert[red]{magneetivuo on kvantisoitunut}.} \appear{Magneettivuo joka läpäisee suprajohtavan virtasilmukan (vorteksitila)}\appear{, täytyy olla $\pi \hbar / e$:n monikerta}

\end{itemize}

\item Suprajohtavan aineen tilatiheyden energiarako voidaan arvioida 

\item Esimerkiksi cadmiumille $T_{c}=0,517 \mathrm{~K}$ jolloin:
\[
E_{g}  \equal \appear{3,5 \,k_{B} T_{c}} \equal \appear{1,56 \times 10^{-4} \;\mathrm{eV}}
\]
\end{itemize}

\endofslide 
%\endofsection
\end{frame}
%%%%%%%%%%%%%%%

%%%%%%%%%%%%%%%
\begin{frame}
\frametitle{\texorpdfstring{\culine[\sectiontitleulinecolor]{\sectiontitle}}{\sectiontitle} }
\framesubtitle{Korkean $T_c$:n suprajohteet}%High-$T_c$ superconductors
\appear{}
%$\underline{\text { High-T }_{c} \text { Superconductors }}$
\begin{itemize}
	\item Suprajohteiden tutkimuksessa keskitytään korkean lämpötilan suprajohteisiin\appear{, graalin maljana suprajohtavuus huoneenlämmössä} 
	
	\item Lupaavia yhdisteitä ovat\appear{ yllättäen \CDAlert[red]{keraamit} jotka sisältävät esim. kuparia ja happea }
	
	\item Nestemäisen typen kiehumispisteen saavuttaminen \appear{($T=77 \mathrm{~K}=-196^{\circ} \mathrm{C}$)} \appear{tapahtui historiallisesti tärkeillä yhdisteillä} \appear{$\mathrm{YBa}_{2} \mathrm{Cu}_{3} \mathrm{O}_{7}$} \appear{ja $\mathrm{Hg}_{0,8} \mathrm{Tl}_{0,2} \mathrm{Ba}_{2} \mathrm{Ca}_{3} \mathrm{Cu}_{3} \mathrm{O}_{8,33}$}\appear{, ja viimeaikainen ennätys Fm-3m rakenteella  $\mathrm{LaH}_{10}$ ($P \sim 170 \mathrm{\;GPa})$, $T=250 \mathrm{~K}=$ $-23^{\circ} \mathrm{C}$, Drozdov et al., 2018)!!}

\item Huomaa että tämä on vain lämpötilaennätys\appear{, uudet korkean $T$:n suprajohteet \CDAlert[fuchsia]{eivät kestä suuria sähkövirtoja}}

\item \CDAlert[blue]{2D-materiaaleissa} tapahtuva suprajohtavuuteen johtava mekanismi vaikuttaa samanlaiselta kuin korkean $T$:n materiaaleissa\appear{, toivon mukaan johtaen \CDAlert[blue]{uusiin teoreettisiin läpimurtoihin}}
\end{itemize}

\endofslide 
%\endofsection
\end{frame}
%%%%%%%%%%%%%%%

%%%%%%%%%%%%%%%
\begin{frame}
\frametitle{\texorpdfstring{\culine[\sectiontitleulinecolor]{\sectiontitle}}{\sectiontitle} }
\framesubtitle{Korkean $T_c$:n suprajohteet}
\appear{}
\begin{itemize}
	\item Teknologisesti merkittävissä sovelluskohteissa\appear{, on usein tärkeämpää maksimoida tripletti:} \appear{a) \CDAlert[red]{lämpötila}}\appear{, b) \CDAlert[blue]{sähköinen varaustiheys}} \appear{ja c) \CDAlert[fuchsia]{magneettikentän voimakkuus}}

\end{itemize}

\vspace{2.5mm}
\begin{minipage}{6cm}
\begin{itemize}[itemsep=1.6mm]
\item Joitakin käytännön sovelluskohteita korkean lämpötilan suprajohteille löytyy jo voimansiirrosta: \appear{erityisiä kaapeliratkaisuita on olemassa}\appear{, joissa suprajohtavuus ylläpidetään tehokkaalla \Alert[red]{konsent\-ri\-sella jäähdytyksellä} }

\item \Alert[red]{Nestemäisen typen lämpötilojen} saavuttaminen oli tärkeää, ja on johtanut lukuisiin sovelluskohteisiin! 

\end{itemize}
\end{minipage}

\appear{\xdef\loc{south east}
\begin{tikzpicture}[remember picture,overlay] % anchor=south
    \node (A) [yshift=-0.25*\figYshift,xshift=\figXshift, anchor=\loc] at (current page.\loc) {\scriptsize 
\begin{tabular}{|cc|} 
%&\text { High-Temperature Superconductors }\\
%&\begin{array}{c|c}
\hline
\text { \Alert[red]{Elektroniikka} } & \text { \Alert[blue]{Perustutkimus} } \\
\text { Sensorit } & \text { Suurtehomagneetit } \\
\text{Supertietokoneet} & \text { SQUID-sensorit } \\
\text{\!\!Kryoelektroniset komponentit\!\!\!\!\!\!} & \text { Bolometrit } \\
  & \text { Hiukkasilmaisimet } \\
 \text { \Alert[blue]{Lääketieteelliset tekniikat} } &  \\
\text { SQUID-diagnostiikka } & \text { \Alert[red]{Mikroaaltoteknologia} } \\
\text { Magneettiresonanssi- } & \text { Filterit } \\
\text { kuvantaminen (MRI) } & \text { Resonaattorit } \\
& \text { Antennit } \\
\text { \Alert[red]{Tehoelektroniikka} } &  \text { Viivelinjat}\\
\text { Generaattorit } & \text { Aktiiviset komponentit } \\
%\text { Motors } &  \\
\text { Muuntajat } & \\
\text { Virranrajoittimet } & \text { \Alert[blue]{Liikennetekniikka} }  \\
\text { Kytkimet } & \text { Maglev-junat} \\
\text { Kaapelit } &  \text {Maglev-tekniikka\!\!} \\
\text { Energian varastointi } & \text { Laivojen koneet }  \\
\hline
%\end{array}
\end{tabular}
}; %scale=\figscale. % 8cm max height
\end{tikzpicture}}

%\endofslide 
\endofsection
\end{frame}
%%%%%%%%%%%%%%%
